\section{Deciding abelian-group theory}
\label{sec:p03-abelian-decision}

Szmielew's theorem decides the common first-order theory of abelian groups,
not the complete theory of an arbitrary individual group. The result is
reported in \cite[p.~611, Theorem~7]{Rabin1977Decidable}. Here the module
normal form of \cref{thm:p03-w03-p03-073-pp-elimination} is combined with
integer linear algebra and a finite realizability test to give the
decision procedure and its proof. Choice is used for vector-space bases
and bounded-group decompositions, as in the preceding module theorem.
No group presentation or oracle for invariant cardinals is an algorithm input.

\begin{hbdefinition}{def:p03-ab-setting}{Language and notation}
Let \(\mathsf{Ab}\) be the finite theory in \(\{+,0\}\) consisting of the
universal group-associativity, commutativity and zero identities, together
with \(\forall x\,\exists y\,(x+y=0)\). Integer multiples and differences
are definitional abbreviations. In a linear equation, negative summands
can be moved to the other side, giving an equation in \(\{+,0\}\).
Every model has its unique \(\mathbb Z\)-module expansion, as in
\cref{def:p03-w03-p03-073-language}.

For an abelian group \(G\), write \(mG=\{mx:x\in G\}\) and
\(G[m]=\{x\in G:mx=0\}\). Here integers are external coefficients,
not quantified group elements. The formula \(d\mid x\) means
\(\exists y\,(dy=x)\), where \(y\) is a group element.
The notation \(C_{p^s}\) denotes the additive cyclic group
\(\mathbb Z/p^s\mathbb Z\), for a prime \(p\) and \(s\geq1\).
A superscript \((\kappa)\) denotes a direct sum of \(\kappa\) copies.
\end{hbdefinition}

\subsection{From sentences to unary indices}

\begin{hblemma}{thm:p03-ab-index-normal}{proved=annex:p3.ab.1}{Effective index reduction}
There is an algorithm which, from a sentence of \(\{+,0\}\), produces
a Boolean combination of index-threshold sentences
\[
 [H(G):K(G)]\geq k,
 \qquad k\geq1,
\]
equivalent to it in every abelian group. Each \(H\) and \(K\) is a unary
pp formula, and \(K(G)\subseteq H(G)\) in every group. These are finite
first-order sentence abbreviations in the sense of
\cref{def:p03-w03-p03-073-invariants}.
\end{hblemma}

\begin{hbdefinition}{def:p03-ab-linear-form}{Unary integer form}
A unary integer form is
\[
 (mx=0)\ \land\ \bigwedge_{i=1}^{t}(d_i\mid c_i x),
 \qquad
 m\geq0,\quad d_i\geq1,\quad c_i\in\mathbb Z\setminus\{0\}.
\]
The first conjunct is vacuous when \(m=0\), and the list of divisibility
conditions may be empty. Clauses \(d\mid0x\) are discarded as tautologies.
All coefficients are supplied as finite integer data.
\end{hbdefinition}

\begin{hblemma}{thm:p03-ab-unary-normal}{proved=annex:p3.ab.2}{Integer reduction}
A unary pp formula can be effectively converted to a unary integer form
equivalent in every abelian group. The construction uses integer
unimodular row and column operations; it does not assume that the group
is a vector space or is torsion-free.
\end{hblemma}

Annihilation conditions are retained separately. Nonzero-divisor
conditions \(d\mid cx\) alone all hold throughout the additive group
\(\mathbb Q\), whereas \(px=0\) does not hold at a nonzero rational.
Replacing that annihilation condition by \(x=0\) would also be wrong
uniformly: in \(C_p\), \(px=0\) holds for every element.

\subsection{The finite invariant data}

\begin{hbdefinition}{def:p03-ab-invariants}{Height, torsion and divisibility data}
For a prime \(p\) and \(n\geq0\), put
\[
 U_{p,n}(G)=(p^nG)[p]=p^nG\cap G[p].
\]
The following are dimensions of vector spaces over \(\mathbb F_p\):
\[
 \begin{aligned}
 \alpha_{p,n}(G)&=\dim_{\mathbb F_p}
                   \bigl(U_{p,n}(G)/U_{p,n+1}(G)\bigr),\\
 \beta_{p,n}(G)&=\dim_{\mathbb F_p}U_{p,n}(G),\\
 \gamma_{p,n}(G)&=\dim_{\mathbb F_p}
                   \bigl(p^nG/p^{n+1}G\bigr).
 \end{aligned}
\]
They may be infinite. Only comparisons with finite integers occur in the
algorithm. For example, \(\alpha_{p,n}\geq r\) says that the displayed
quotient has at least \(p^r\) distinct cosets; it is a finite first-order
sentence. The analogous descriptions apply to \(\beta\) and \(\gamma\).
For \(m\geq1\), let
\[
 \operatorname{Ann}_m:\quad \forall x\,(mx=0).
\]
The indices \(p,n,r,m\) are metatheoretic integers, not additional
variables or sorts in the group language. For a cardinal \(\kappa\)
and a positive integer \(b\), define
\[
 \operatorname{cap}_b(\kappa)=
 \begin{cases}\kappa,&\kappa<b,\\b,&\kappa\geq b.\end{cases}
\]
Its values are the integers \(0,\ldots,b\).
\end{hbdefinition}

\begin{hblemma}{thm:p03-ab-bounded}{proved=annex:p3.ab.3}{Bounded groups}
If \(p^E G=0\), then
\[
 G\cong\bigoplus_{s=1}^{E}
             C_{p^s}^{(\alpha_{p,s-1}(G))}.
\]
A group annihilated by a positive integer is the direct sum of its
finitely many primary components, to which this statement applies.
The multiplicities can be infinite; bounded exponent is not a
finiteness assertion.
\end{hblemma}

\begin{hblemma}{thm:p03-ab-layers}{proved=annex:p3.ab.4}{Layer comparisons}
For every abelian group,
\[
 \beta_{p,n}=\alpha_{p,n}+\beta_{p,n+1},
 \qquad
 \gamma_{p,n}=\alpha_{p,n}+\gamma_{p,n+1}.
\]
These are cardinal sums, not subtraction identities. For \(E\geq1\),
the cyclic multiplicities in the two bounded groups below are:
\[
 \begin{array}{c|c|c}
 &C_{p^s},\ 1\leq s<E&C_{p^E}\\ \hline
 G[p^E]&\alpha_{p,s-1}(G)&\beta_{p,E-1}(G)\\
 G/p^EG&\alpha_{p,s-1}(G)&\gamma_{p,E-1}(G).
 \end{array}
\]
\end{hblemma}

\begin{hbdefinition}{def:p03-ab-cyclic-data}{Cyclic solution data}
For a unary integer form \(\varphi\), let
\[
 a_\varphi(p,s)=
 \begin{cases}
 0,&m=0,\\
 \max(s-v_p(m),0),&m>0,
 \end{cases}
\]
where \(v_p\) is the exponent of \(p\) in a nonzero integer, using its
absolute value. Define
\[
 h_\varphi(p,s)=
 \max\bigl(\{0,a_\varphi(p,s)\}
       \cup\{\min(s,v_p(d_i))-v_p(c_i):1\leq i\leq t\}\bigr).
\]
For \(K\subseteq H\), put
\[
 w_{H,K}(p,s)=h_K(p,s)-h_H(p,s).
\]
The solution subgroup in \(C_{p^s}\) is \(p^{h_\varphi(p,s)}C_{p^s}\);
hence \(w_{H,K}(p,s)\geq0\) and the index on this cyclic group is
\(p^{w_{H,K}(p,s)}\). These assertions are verified in the next proof.
\end{hbdefinition}

\begin{hbtheorem}{thm:p03-ab-index-calculation}{proved=annex:p3.ab.5}{Computing invariant conditions}
Every unary pp-index threshold can be effectively expressed as a Boolean
combination of sentences
\[
 \alpha_{p,n}\geq r,\qquad
 \beta_{p,n}\geq r,\qquad
 \gamma_{p,n}\geq r,\qquad
 \operatorname{Ann}_m.
\]
The equivalence holds in every abelian group, including uncountable
groups and groups of unbounded exponent.
\end{hbtheorem}

\begin{hbmethod}{def:p03-ab-index-algorithm}{The three index calculations}
For an index threshold \(k=1\), output truth. Otherwise replace \(K\)
by \(H\land K\), compute the two unary integer forms, and let \(D\)
be a positive common multiple of all their divisors \(d_i\), taking
\(D=1\) for an empty list.
\begin{enumerate}
 \item If \(m_H=m_K=0\), compute in \(G/DG\). For each \(p^E\) exactly
       dividing \(D\), use the quotient row of
       \cref{thm:p03-ab-layers}.
 \item If \(m_H>0\), put \(L=m_HD\) and compute in \(G[L]\). For each
       \(p^E\) exactly dividing \(L\), use the torsion-subgroup row.
 \item If \(m_H=0<m_K\), put \(L=m_KD(k-1)!\).
       Emit truth on the branch \(\neg\operatorname{Ann}_L\).
       On the branch \(\operatorname{Ann}_L\), use the bounded-group
       multiplicities \(\alpha_{p,s-1}(G)\), for \(p^s\mid L\).
\end{enumerate}
In each bounded calculation, replace every relevant multiplicity by
a value \(0,\ldots,k\), with \(k\) meaning at least \(k\). Test the finite
integer product
\[
 \prod_{p,s}p^{\,w_{H,K}(p,s)\operatorname{cap}_k(\kappa_{p,s})}
       \ \geq\ k.
\]
Enumerate the finitely many capped assignments for which it holds,
and express each assignment by finite dimension-threshold conditions.
These are symbolic calculations producing a sentence. The procedure
does not inspect an input group or compute its actual infinite cardinals.
\end{hbmethod}

\subsection{Finite realizability and the decision procedure}

\begin{hbdefinition}{def:p03-ab-test-groups}{Explicit comparison groups}
The Prüfer \(p\)-group and the \(p\)-local rational group are
\[
 C_{p^\infty}=\mathbb Z[1/p]/\mathbb Z,\qquad
 \mathbb Z_{(p)}=\{a/q\in\mathbb Q:p\nmid q\},
\]
with their additive group operations and nonzero denominators.
For fixed finite data and a positive truncation bound \(b\), the
comparison groups are:
\[
 \bigoplus_{\substack{p\mid d\\1\leq s\leq v_p(d)}}
                 C_{p^s}^{(a_{p,s})}
 \quad (0\leq a_{p,s}\leq b),
\]
and
\[
 \mathbb Q\ \oplus
 \bigoplus_{p\in P}
 \left(
   \bigoplus_{j=0}^{N_p}C_{p^{j+1}}^{(a_{p,j})}
   \ \oplus C_{p^\infty}^{(b_p)}
   \ \oplus \mathbb Z_{(p)}^{(c_p)}
 \right),
\]
where every multiplicity lies in \(0,\ldots,b\).
These are finite descriptions of actual groups; the second family is
not a claim that arbitrary abelian groups have such a decomposition.
\end{hbdefinition}

\begin{hbtheorem}{thm:p03-ab-realization}{proved=annex:p3.ab.6}{Finite realizability}
Satisfiability of any finite Boolean combination of the sentences in
\cref{thm:p03-ab-index-calculation} is decidable. When it is satisfiable,
the procedure returns one of the explicitly described countable groups
in \cref{def:p03-ab-test-groups} satisfying it.
\end{hbtheorem}

\begin{hbmethod}{def:p03-ab-realization-algorithm}{Testing the finite conditions}
Put the Boolean combination in disjunctive normal form and test each
conjunction. Let \(b\geq1\) be the largest finite dimension threshold
in it, or one if there are none.
\begin{enumerate}
 \item If the conjunction contains positive annihilator sentences,
       let \(d\) be the greatest common divisor of their indices.
       Enumerate the first family of comparison groups and test the
       literals by their finite cyclic multiplicities.
 \item Otherwise let \(P\) be the finite set of primes mentioned in
       dimension conditions, and \(N_p\) the largest measured level
       for \(p\). Enumerate the second family. The rational summand
       satisfies every negative annihilator sentence. At the measured
       levels its dimension data are
       \[
       \begin{aligned}
       \alpha_{p,n}&=a_{p,n},\\
       \beta_{p,n}&=\sum_{j=n}^{N_p}a_{p,j}+b_p,\\
       \gamma_{p,n}&=\sum_{j=n}^{N_p}a_{p,j}+c_p.
       \end{aligned}
       \]
       Components indexed by other primes contribute zero.
\end{enumerate}
A successful finite test returns its group description. If no candidate
for any conjunction succeeds, report unsatisfiability. The proof shows
why the computed bounds cover all possible models, rather than merely
a selected finite collection of examples.
\end{hbmethod}

\begin{hbtheorem}{thm:p03-abelian-decidable}{proved=annex:p3.ab.7}{Szmielew's theorem}
There are algorithms deciding satisfiability and validity of first-order
sentences in \(\{+,0\}\) relative to the class of all abelian groups.
Every input sentence has an effectively obtainable equivalent Boolean
combination of the explicit finite invariant conditions above.
For a sentence which is not valid, a finite description of a countable
countermodel is effectively obtainable.
\end{hbtheorem}

\begin{hblemma}{thm:p03-ab-individual-boundary}{proved=annex:p3.ab.8}{An individual theory}
There is a countable abelian group whose complete first-order theory
is undecidable. Thus the common-theory decision procedure does not
supply a uniform decision of the complete theory of each individual group.
\end{hblemma}
